\documentclass[12pt]{article}
\usepackage[utf8]{inputenc}
\usepackage[T1]{fontenc}
\usepackage[margin=1in]{geometry}
\usepackage{setspace}
\usepackage{xcolor}
\usepackage[colorlinks=true, urlcolor=blue!60!black, linkcolor=blue!60!black]{hyperref}
\setlength{\parindent}{2em}
\setlength{\parskip}{1.25em}
\renewcommand{\baselinestretch}{1.0}
\begin{document}
\begin{center}
{\LARGE \textit{Counterman v. Colorado}, 600 U.S. 66 (2023)\par}
\vspace{0.8em}
{\large Case Brief --- Four-Page Edition\par}
\end{center}
\vspace{1em}

\section{Memory Jogger}

A man convicted of stalking a singer through hundreds of Facebook messages obtained vacatur and remand by a 7--2 vote: a five-Justice majority held that a true-threats prosecution requires proof that the speaker was at least reckless---that he consciously disregarded a substantial risk his words would be viewed as threatening violence---not merely that a reasonable person would see a threat (\href{https://supreme.justia.com/cases/federal/us/600/22-138/}{600 U.S. 66}).

\section{Detailed Case Facts}

From 2014 to 2016, Billy Raymond Counterman sent hundreds of unsolicited Facebook messages to C.W., Colorado singer-songwriter Coles Whalen. They had never met, she never replied, and each time she blocked him he opened a new account. Some messages were mundane (``Good morning sweetheart''), some implied surveillance (``Was that you in the white Jeep?''), and several envisaged violent harm: ``Fuck off permanently''; ``Staying in cyber life is going to kill you''; and ``You're not being good for human relations. Die.'' Whalen suffered severe anxiety and insomnia, stopped walking alone, and canceled performances before contacting authorities.

\section{Procedural History}

Colorado charged Counterman under a stalking statute criminalizing repeated communication that would cause a reasonable person serious emotional distress and actually does so, relying solely on the messages. Applying Colorado's objective ``reasonable person'' test for true threats, the trial court rejected his First Amendment motion; a jury convicted, and he was sentenced to four and a half years. The Colorado Court of Appeals affirmed, and the Supreme Court granted certiorari to resolve a split left open after \textit{Elonis v. United States}. It vacated and remanded on June 27, 2023.

\section{Judicial Votes}

Vacated and remanded, \textbf{7--2}.
\begin{itemize}
\item \textbf{Majority:} Justice Kagan, joined by Chief Justice Roberts and Justices Alito, Kavanaugh, and Jackson.
\item \textbf{Concurring in part and in the judgment:} Justice Sotomayor, joined by Justice Gorsuch as to Parts I, II, III-A, and III-B.
\item \textbf{Dissents:} Justice Thomas; Justice Barrett, joined by Justice Thomas.
\end{itemize}

\section{Holding}

True threats of violence are unprotected speech defined objectively by what the statement conveys, but the First Amendment requires the State in a true-threats prosecution to prove that the defendant had at least a reckless mental state, meaning he ``consciously disregarded a substantial risk that his communications would be viewed as threatening violence.'' Because Colorado prosecuted Counterman under an objective standard without proof of subjective awareness, his conviction violated the First Amendment.

\section{Analysis of Opinions}

\textbf{Majority (Kagan).} True threats---serious expressions of intent to commit unlawful violence, as in \textit{Virginia v. Black}---are unprotected, unlike jests or the political hyperbole of \textit{Watts}. Yet the First Amendment often gives ``strategic protection,'' a buffer shielding some unprotected speech to avoid chilling protected expression, as in defamation (\textit{New York Times v. Sullivan}'s actual malice), incitement (\textit{Brandenburg}'s specific intent), and obscenity (scienter). An objective standard forces speakers to give threats ``a wide berth.'' Recklessness strikes the right balance: higher standards would impede prosecutions of culpable speakers; recklessness has given ample ``breathing space'' in defamation; and incitement demands more only because it sits ``a hair's-breadth away'' from political advocacy.

\textbf{Sotomayor (concurring in part).} Recklessness suffices for stalking, which involves repeated targeted contact, but the Court should not have reached pure threats without briefing; writing alone in part, she would require purpose or knowledge to avoid criminalizing online hyperbole, rap lyrics, and protest and to keep prosecutors from evading \textit{Brandenburg}.

\textbf{Thomas (dissenting).} The majority's ``surprising and misplaced reliance'' on \textit{Sullivan}, a ``policy-driven'' departure from common law, compounds error.

\textbf{Barrett (dissenting).} Most unprotected categories use objective standards, and contextual review under \textit{Watts} and \textit{Black} already prevents chill. The recklessness floor is a legislative ``Goldilocks judgment'' that will obstruct civil protective orders, civil enforcement, and school or workplace discipline against delusional threat-makers.

\section{Examples: Future Applications}

\begin{enumerate}
\item \textbf{Same-Side (Conviction Invalid): Zoning Livestream.} A resident posts that ``someone ought to burn his house down tonight'' about a councilman, and the jury is told to convict if a ``reasonable viewer'' would see a threat. The objective instruction fails; on remand, \textit{Watts} hyperbole and \textit{Brandenburg} also matter.
\item \textbf{Same-Side (No Recklessness): Delusional Email.} A reader in a paranoid delusion sends one email, ``your final chapter ends Friday,'' believing it a harmless prophecy. Without awareness of the risk, conviction fails.
\item \textbf{Opposite-Side (Recklessness Provable): Messages After Notice.} The same reader sends 150 messages after ten blocked accounts and a letter saying ``Your messages terrify me.'' Continuing after notice supports recklessness.
\item \textbf{Opposite-Side (Conviction Stands): Intent Statute.} A man texts a photo of his former partner's house saying she will not live to see the lawsuit through, and a jury finds he meant to cause fear of death. Recklessness is the minimum; a State may require more.
\item \textbf{Fence-Sitter: Civil Protective Orders.} A protective order issued under an objective standard restricts a coworker's contact. \textit{Counterman}'s chill rationale reaches any sanction, but the case was criminal, and Barrett warned of civil effects.
\end{enumerate}

\section{Critique}

\textbf{The delusional stalker paradox.} As Justice Barrett noted, genuinely delusional stalkers like Counterman, who appeared to believe he was in an ongoing dialogue with Whalen, may pose the gravest danger to victims while lacking subjective awareness that their words are frightening.

\textbf{Tension with history and tradition.} Unlike \textit{Bruen} or \textit{Stevens}, the majority openly balanced interests, ``having much of what is important on both sides of the scale,'' to build a prophylactic buffer, drawing fire from originalists (Thomas and Barrett) and from speech-protective Justices (Sotomayor and Gorsuch).

\section{Key Quotations}

\begin{quote}
``We hold that it does, but that a mental state of recklessness is sufficient. The State must show that the defendant consciously disregarded a substantial risk that his communications would be viewed as threatening violence. The State need not prove any more demanding form of subjective intent to threaten another.''
\end{quote}
\noindent --- Justice Kagan, Opinion of the Court, 600 U.S. at 69.

\begin{quote}
``Yet the First Amendment may still demand a subjective mental-state requirement shielding some true threats from liability. The reason relates to what is often called a chilling effect. Prohibitions on speech have the potential to chill, or deter, speech outside their boundaries.''
\end{quote}
\noindent --- Justice Kagan, Opinion of the Court, 600 U.S. at 75.

\begin{quote}
``The reality is that recklessness is not grounded in law, but in a Goldilocks judgment: Recklessness is not too much, not too little, but instead `just right.'\thinspace''
\end{quote}
\noindent --- Justice Barrett, dissenting, 600 U.S. at 119.

\end{document}
